\documentclass[../main.tex]{subfiles}


\begin{document}

\setcounter{section}{6}
\externaldocument{../Chapter/C6A.tex}

\begin{exercise}
  Prove or disprove: If $\join{v}{n - 1} \in V$, then:
  \[
  \sum_{j = 0}^{n - 1}\sum_{k = 0}^{n - 1} \ip{v_j}{v_k} \ge 0
  \]
\end{exercise}
\begin{proof}
  For any $j$, we have $\sum_{k = 0}^{n - 1} \ip{v_j}{v_k} = \ip{v_j}{\join[+]{v}{n - 1}}$,
  thus the equation is now $\ip{\join[+]{v}{n - 1}}{\join[+]{v}{n - 1}} \ge 0$,
  which is trivial by definition.
\end{proof}

\begin{exercise}
  Let $S \in \LT(V)$. Define $\ip{\cdot}{\cdot}_1$ by
  \[
  \ip{u}{v}_1 = \ip{Su}{Sv}
  \]
  for all $u, v \in V$.
  Show that $\ip{\cdot}{\cdot}_1$ is an inner product over $V$
  $\iff$ $S$ is injective.
\end{exercise}
\begin{proof}
  ~
  \begin{itemize}
    \item Try to prove $\langle \cdot, \cdot \rangle_1$ is a inner product and see where we stuck without injective.
          We find that $\langle v, v \rangle_1 = 0 \iff v = 0$ stuck, as if
          $S$ is not injective, then we have $\ip{u - v}{u - v}_1 = \ip{S(u - v)}{S(u - v)} = \ip{0}{0} = 0$
          where $u - v \neq 0$ and $Su = Sv$.
          Therefore $S$ must be injective.

          Although the proof is not *constructive*, it help us to find a constructive proof:
          Suppose $Su = Sv$, then $\ip{u - v}{u - v}_1 = \ip{Su - Sv}{Su - Sv} = 0$, thus $u - v = 0$, therefore $u = v$.
    \item Positivity, Additivity, Homogeneity holds cause $S$ is a linear map and $\ip{\cdot}{\cdot}$ is an inner product,
          and Conjugate Symmetry holds cause $\ip{\cdot}{\cdot}$ is an inner product.
          For Definiteness, $\ip{v}{v}_1 = \ip{Sv}{Sv} = 0$ implies $v = 0$ cause $S$ injective,
          and $\ip{0}{0}_1 = \ip{S0}{S0} = \ip{0}{0} = 0$.
  \end{itemize}
\end{proof}

\begin{exercise}
  ~
  \begin{itemize}
    \item Show that $f((a, b), (c, d)) = |ac| + |bd|$ is not an inner product over $\mathbb{R}^2$.
    \item Show that $f((a, b, c), (x, y, z)) = ax + cz$ is not an inner product over $\mathbb{R}^3$.
  \end{itemize}
\end{exercise}
\begin{proof}
  ~
  \begin{itemize}
    \item $f((1, 1), (1, 1)) = 1 + 1$ and $f((-1, -1), (1, 1)) = 1 + 1$, then $f((1, 1) + (-1, -1), (1, 1)) = f((0, 0), (1, 1)) = 0 + 0 = 0$
          while $f((1, 1) + (-1, -1), (1, 1)) = f((1, 1), (1, 1)) + f((-1, -1), (1, 1)) = 2 + 2 = 4$.
    \item $f((0, 1, 0), (0, 1, 0)) = 0$ but $(0, 1, 0) \neq 0$.
  \end{itemize}
\end{proof}

\begin{exercise}
  Let $T \in \LT(V)$ and $\norm{Tv} \le \norm{v}$ for all $v \in V$.
  Show that $T - \sqrt{2}I$ is injective.
\end{exercise}
\begin{proof}
  Suppose $T - \sqrt{2}$ is not injective, then $Tv = \sqrt{2}v$ for some $v$,
  then $\norm{Tv} = \norm{\sqrt{2}v} = |\sqrt{2}|\norm{v} \ge \norm{v}$.
  Basically any eigenvalue with absolute value greater than $1$ can make it.
\end{proof}

\begin{exercise}
  Let $V$ an inner product space over $\mathbb{R}$.
  \begin{itemize}
    \item Show that $\ip{u + v}{u - v} = \norm{u}^2 - \norm{v}^2$.
    \item Show that $u + v \perp u - v$ if $\norm{u} = \norm{v}$.
    \item Use last conclusion to show that the diagonal of 菱形 are orthogonal.
  \end{itemize}
\end{exercise}
\begin{proof}
  ~
  \begin{itemize}
    \item $\ip{u + v}{u - v} = \ip{u}{u} + \ip{v}{u} - \ip{u}{v} - \ip{v}{v}$,
          note that $\ip{v}{u} = \ip{u}{v}$ since $V$ over $\mathbb{R}$,
          thus $\ip{u + v}{u - v} = \ip{u}{u} - \ip{v}{v} = \norm{u}^2 - \norm{v}^2$.
    \item By $\Uparrow$.
    \item We know 菱形 is a parallelogram that four sides have same length,
          thus $\norm{u} = \norm{v}$ and $u + v \perp u - v$, where $u + v$ and $u - v$
          are two diagonal of 菱形.
  \end{itemize}
\end{proof}

\begin{exercise}
  Let $u, v \in V$. Show that $\ip{u}{v} = 0$ $\iff$ $\norm{u} \le \norm{u + av}$ for any $a \in F$.
\end{exercise}
\begin{proof}
  ~
  \begin{itemize}
    \item $(\Rightarrow)$ We will show that $\norm{u} \le \norm{u + av}$ by $\norm{u}^2 \le \norm{u + av}^2$
          (recall that norm is always non-negative).
          $\norm{u + av}^2 = \norm{u}^2 + (a\norm{v})^2$
    \item $(\Leftarrow)$ If $v = 0$, then $0 \perp u$ and the proof is complete, we assume $v \neq 0$. Let $w \in V$ such that $cv + w = u$ and $v \perp w$,
          then $\norm{u}^2 = \norm{cv + w}^2 = \norm{cv}^2 + \norm{w}^2$ (see \cref{T6.13}), thus $\norm{u}^2 \ge \norm{w}^2$,
          therefore $\norm{u} \ge \norm{w}$ where $w = u - cv$, therefore $\norm{w} \le \norm{w}$,
          hence $\norm{u} = \norm{w}$.
          Then $\norm{u}^2 = \norm{cv}^2 = \norm{w}^2$ is now $\norm{cv}^2 = 0$, therefore $c = 0$ or $v = 0$, but $v \neq 0$,
          thus $c = \displaystyle \frac{\ip{u}{v}}{\norm{v}^2} = 0$ then $\ip{u}{v} = 0$ and $u \perp v$.
  \end{itemize}
\end{proof}

\begin{exercise}
  Let $u, v \in V$. Show that $\norm{au + bv} = \norm{bu + av}$ for any $a, b \in \mathbb{R}$ $\iff$ $\norm{u} = \norm{v}$.
\end{exercise}
\begin{proof}
  ~
  \begin{itemize}
    \item We have $\norm{au + bv}^2 = \ip{au}{au} + \ip{bv}{bv} + \ip{au}{bv} + \ip{bv}{au}$
          and     $\norm{bu + av}^2 = \ip{bu}{bu} + \ip{av}{av} + \ip{bu}{av} + \ip{av}{bu}$.
          Note that $\ip{au}{bv} = a\bar{b}\ip{u}{v} = \bar{a}b\ip{u}{v} = \ip{bu}{av}$,
          similarly, $\ip{bv}{au} = \ip{av}{bu}$.
          Thus $\ip{au}{au} + \ip{bv}{bv} = \ip{bu}{bu} + \ip{av}{av}$, let $a = \sqrt{2}$ and $b = 1$,
          we have $\sqrt{2}^2\ip{u}{u} + \ip{v}{v} = \ip{u}{u} + \sqrt{2}^2\ip{v}{v}$
          therefore $\ip{u}{u} = \ip{v}{v}$, thus $\norm{u} = \norm{v}$.
    \item $\norm{au + bv}^2 - \norm{bu + av}^2 = a^2\ip{u}{u} + b^2\ip{v}{v} + b^2\ip{u}{u} + a^2\ip{a}{a} = 0$.
  \end{itemize}
\end{proof}

\setcounter{exercise}{8}
\begin{exercise}
  Let $\norm{u} = \norm{v} = 1$ and $\ip{u}{v} = 1$. Show that $u = v$.
\end{exercise}
\begin{proof}
  We know $\ip{u}{v} = \ip{v}{u} = 1$ and then
  $\ip{u}{v} - \ip{v}{v} = 1 - 1 = 0 = \ip{u - v}{v}$ and
  $- \ip{v}{u} + \ip{u}{u} = - 1 + 1 = 0 = \ip{u - v}{u}$ therefore 
  $\ip{u - v}{v} - \ip{u - v}{u} = \ip{u - v}{u - v} = 0$,
  thus $u - v = 0$.
\end{proof}

\begin{exercise}
  Let $u, v \in V$ and $\norm{u}, \norm{v} \le 1$. Show that:
  \[
  \sqrt{1 - \norm{u}^2}\sqrt{1 - \norm{v}^2} \le 1 - |\ip{u}{v}|
  \]
\end{exercise}
\begin{proof}
  By Cauchy-Schwarz inequality, we have $1 - |\ip{u}{v}| \ge 1 - \norm{u}\norm{v}$.
  We will show that $\sqrt{1 - \norm{u}^2}\sqrt{1 - \norm{v}^2} \le 1 - \norm{u}\norm{v}$.
  \begin{align*}
    \sqrt{1 - \norm{u}^2}\sqrt{1 - \norm{v}^2} & \le 1 - \norm{u}\norm{v} \\
    (1 - \norm{u}^2)(1 - \norm{v}^2) & \le (1 - \norm{u}\norm{v})^2 \\
    1 - \norm{u}^2 - \norm{v}^2 + \norm{u}^2\norm{v}^2 & \le 1 - 2\norm{u}\norm{v} + \norm{u}^2\norm{v}^2 \\
    - \norm{u}^2 - \norm{v}^2 & \le - 2\norm{u}\norm{v} \\
    \norm{u}^2 + \norm{v}^2 & \ge 2\norm{u}\norm{v} \\
    \norm{u}^2 - 2\norm{u}\norm{v} + \norm{v}^2 & \ge 0 \\
    (\norm{u} - \norm{v})^2 & \ge 0 \\
  \end{align*}
  The second step is well since both sides are non-negaive.
\end{proof}

\begin{exercise}
  Find $u, v \in \mathbb{R}^2$, such that $u$ a scalar multiple of $(1, 3)$
  and $v$ is perpenticular to $(1, 3)$ and $(1, 2) = u + v$.
\end{exercise}
\begin{proof}
  人话：找到一个关于 $(1, 3)$ 和垂直于 $(1, 3)$ 的向量 的线性组合，使得它等于 $(1, 2)$.

  We have $(3, -1) \perp (1, 3)$, then $x = \frac{7}{10}$ and $y = \frac{1}{10}$
  is a solution of:
  \begin{align*}
    x + 3y &= 1 \\
    3x - 1y &= 2 \\
  \end{align*}
  Thus $u = \frac{7}{10}(1, 3)$ and $v = \frac{1}{10}(3, -1)$.
\end{proof}

\setcounter{exercise}{13}
\begin{exercise}
  Let $v \in V$ and $v \neq 0$.
  Show that for any $u \in V$ and $\norm{u} = 1$, then
  \[
  \norm{v - \frac{v}{\norm{v}}} \le \norm{v - u}
  \]
  equal only if $u = \frac{v}{\norm{v}}$.

  这段话说的是：$v$ 穿过单位圆的点的向量是距离 $v$ 最近的元素。
\end{exercise}
\begin{proof}
  Let $v^\prime = \frac{v}{\norm{v}}$ and $w \in V$ where $w \perp v^\prime$ such that $u = cv^\prime + w$,
  we can see $|c| \le 1$ by 勾股定理. Then
  \begin{align*}
    \norm{v - v^\prime} & \le \norm{v - u} \\
    \norm{v - v^\prime}^2 & \le \norm{v - u}^2 \\
    \ip{v}{v} + \ip{-v^\prime}{-v^\prime} + \ip{v}{-v^\prime} + \ip{-v^\prime}{v} & \le \ip{v}{v} + \ip{-u}{-u} + \ip{v}{-u} + \ip{-u}{v} \\
    \ip{v}{-v^\prime} + \ip{-v^\prime}{v} & \le \ip{v}{-u} + \ip{-u}{v} \\
    - \ip{v}{v^\prime} - \ip{v^\prime}{v} & \le - \ip{v}{u} - \ip{u}{v} \\
    (-1) (\ip{v}{v^\prime} + \overline{\ip{v}{v^\prime}}) & \le (-1) (\ip{v}{cv^\prime + w} + \ip{cv^\prime + w}{v}) \\
    \ip{v}{v^\prime} + \overline{\ip{v}{v^\prime}} & \ge \ip{v}{cv^\prime + w} + \ip{cv^\prime + w}{v} \\
    \ip{v}{v^\prime} + \ip{v}{v^\prime} & \ge \ip{v}{cv^\prime} + \ip{v}{w} + \ip{cv^\prime}{v} + \ip{w}{v} \\
    2 \ip{v}{v^\prime} & \ge \ip{v}{cv^\prime} + \ip{cv^\prime}{v} \\
    2 \ip{v}{v^\prime} & \ge \bar{c} \ip{v}{v^\prime} + c \ip{v^\prime}{v} \\
    2 \ip{v}{v^\prime} & \ge \bar{c} \ip{v}{v^\prime} + c \overline{\ip{v}{v^\prime}} \\
    2 \ip{v}{v^\prime} & \ge \bar{c} \ip{v}{v^\prime} + c \ip{v}{v^\prime} \\
    2 & \ge \bar{c} + c \\
    2 & \ge 2 \real c \\
    1 & \ge \real c \\
  \end{align*}
  We have $\real c \le |c|$ and $|c| \le 1$, thus $\real c \le 1$.

  The inequality equals if $\real c = 1$, which means $c = 1$ as $|c| = 1$ so the imagine part is $0$,
  thus $u = cv^\prime + w = v^\prime + w$ while $\norm{u} = \norm{v^\prime} = 1$, which means $w = 0$,
  therefore $u = v^\prime$.
\end{proof}

\begin{exercise}
  Let non-zero $u, v \in \mathbb{R}^2$. Show that:
  \[
  \ip{u}{v} = \norm{u} \norm{v} \cos \theta
  \]
  where $\theta$ is the angle between $u$ and $v$.
\end{exercise}
\begin{proof}
  Note that the inner product of $\mathbb{R}^2$ is euclidian inner product,
  then let $u = a (\cos \alpha, \sin \alpha)$ and $v = b (\cos \beta, \sin \beta)$,
  thus
  \begin{align*}
    \ip{u}{v} &= ab (\cos \alpha \cos \beta) + ab (\sin \alpha \sin \beta) \\
              &= ab (\cos \alpha \cos \beta + \sin \alpha \sin \beta) \\
              &= ab (\cos \alpha \cos (- \beta) - \sin \alpha \sin (- \beta)) \\
              &= ab (\cos (\alpha - \beta)) \\
              &= \norm{u} \norm{v} (\cos (\alpha - \beta)) \\
  \end{align*}
  Recall that $\norm{u} = a$ and $\norm{v} = b$, since $\norm{(\cos \alpha, \sin \alpha)} = 1$,
  similar to $b$. Also $\cos \theta = \cos - \theta$ and $\sin - \theta = - \sin \theta$.
\end{proof}

\begin{exercise}
  The angle between two vectors in $R^n$ are hard to see when $n > 3$, thus define the angle between $u, v \in R^n$
  by:
  \[
  \arccos \frac{\ip{x}{y}}{\norm{x}\norm{y}}
  \],
  try explain why this requires Cauchy-Schwarz inequality.
\end{exercise}
\begin{proof}
  Note that $\arccos$ is the inverse of $\cos$, thus $\arccos$ only defines
  at $[ -1 , 1 ]$, which means $|\frac{\ip{x}{y}}{\norm{x}\norm{y}}| \le 1$,
  Cauchy-Schwarz inequality says $|\ip{x}{y}| \le \norm{x}\norm{y}$, thus $|\frac{\ip{x}{y}}{\norm{x}\norm{y}}| \le 1$.
\end{proof}

\setcounter{exercise}{19}
\begin{exercise}
  Show that $u, v \in V$, then $|\norm{u} - \norm{v}| \le \norm{u - v}$.
\end{exercise}
\begin{proof}
  We will show $|\norm{u} - \norm{v}|^2 \le \norm{u - v}^2$,
  we can see that no matter $\norm{u} - \norm{v}$ greater or less than $0$,
  $|\norm{u} - \norm{v}|^2$ is always $\norm{u}^2 + \norm{v}^2 - 2\norm{u}\norm{v}$,
  thus
  \begin{align*}
    \norm{u}^2 + \norm{v}^2 - 2\norm{u}\norm{v} & \le \ip{u}{u} + \ip{v}{v} + \ip{u}{- v} + \ip{- v}{u} \\
    - 2\norm{u}\norm{v} & \le (-1) (\ip{u}{v} + \ip{v}{u}) \\
    - 2\norm{u}\norm{v} & \le (-1) (\ip{u}{v} + \overline{\ip{u}{v}}) \\
    - 2\norm{u}\norm{v} & \le (-2) \real \ip{u}{v} \\
    \norm{u}\norm{v} & \ge \real \ip{u}{v} \\
  \end{align*}
  Recall that $|\ip{u}{v}| \le \norm{u}\norm{v}$ and $\real \ip{u}{v} \le |\real \ip{u}{v}| \le |\ip{u}{v}|$.
\end{proof}

\setcounter{exercise}{21}
\begin{exercise}
  Let $u, v \in V$, show that
  \[
  \norm{u + v} \norm{u - v} \le \norm{u}^2 + \norm{v}^2
  \]
\end{exercise}
\begin{proof}
  The the left hand side is:
  \begin{align*}
     & \norm{u + v} \norm{u - v} \\
    =& \sqrt{\ip{u + v}{u + v} \ip{u - v}{u - v}} \\
    =& \sqrt{(\norm{u}^2 + \norm{v}^2 + \ip{u}{v} + \ip{v}{u})(\norm{u}^2 + \norm{v}^2 + \ip{u}{-v} + \ip{-v}{u})} \\
    =& \sqrt{(\norm{u}^2 + \norm{v}^2 + (\ip{u}{v} + \ip{v}{u}))(\norm{u}^2 + \norm{v}^2 - (\ip{u}{v} + \ip{v}{u}))} \\
  \end{align*}
  let $a = \norm{u}^2 + \norm{v}^2$ and $b = \ip{u}{v} + \ip{v}{u}$, then
  \begin{align*}
     & \norm{u + v} \norm{u - v} \\
    =& \sqrt{(a + b)(a - b)} \\
    =& \sqrt{a^2 - b^2} \\
  \end{align*}
  while the right hand side is: $\norm{u}^2 + \norm{v}^2 = a$, then square both sides we have:
  \[
  a^2 - b^2 \le a^2
  \]
  which is obviously true, cause $b^2 \ge 0$ (it is not the case that $b$ is a complex number, as $b = \norm{u + v}^2 - \norm{u}^2 - \norm{v}^2$).
\end{proof}
    
\end{document}